
\documentclass[11pt]{article}
\usepackage[margin=1in]{geometry}
\usepackage{amsmath,amssymb,amsthm}
\usepackage{hyperref}
\usepackage{enumitem}

\newtheorem{remark}{Remark}

\title{Quarantined Computational Material (Non-Proofs): Hessian Spectra, Projectors, HOTRG Jacobians}
\date{}

\begin{document}
\maketitle

\section{Status}

This document collects computational outputs and code-level constructions from the project directory.
Everything here is \emph{evidence only}.  None of it is treated as a proof of any analytic statement.

\section{Finite-volume Hessian spectra at vacuum}

\subsection{Wilson Hessian on $L=2$ (projected physical sector)}

A JAX-based Hessian construction with an explicit gauge+toron projector reports, at $U\equiv I$ and $\beta=1$:
\begin{itemize}[leftmargin=*]
\item $\lambda_{\max}(H_0)\approx 2.6666667=8/3$,
\item $\lambda_{\min}(H_0|\_{\mathrm{phys}})\approx 0.6666667=2/3$.
\end{itemize}

These values are consistent with the analytic identity $H_0=(1/6)\Delta_1$ from Document~02 when the discrete $1$-form Laplacian has smallest positive eigenvalue $4$ on $L=2$.

\subsection{Discrete $1$-form Laplacian on $L=2$}

A separate construction of the coboundary matrix $d$ (links $\to$ plaquettes) and $\Delta_1=d^\top d$ on an $L=2$ periodic $4$D lattice reports:
\begin{itemize}[leftmargin=*]
\item total dimension $|E_L|\cdot \dim(\mathfrak{su}(3)) = 64\cdot 8 = 512$ (per color suppressed in the Laplacian-only test),
\item $\dim\ker(\Delta_1)=1040$ in the scalarized test (counting gauge+toron structure in the chosen representation),
\item smallest strictly positive eigenvalue in the projected sector is $4$.
\end{itemize}

\section{Scaling with $L$ (projected Lanczos at vacuum)}

A projected Lanczos scan at vacuum (with a Fourier-space transverse projector) reports:
\begin{itemize}[leftmargin=*]
\item for $L=2$: $\lambda_{\min}\approx 0.66667$,
\item for $L=3$: $\lambda_{\min}\approx 0.50000$,
\end{itemize}
again consistent with Document~02, where $\lambda_{\min}(H_0|\_{\mathrm{phys}})=\frac{1}{6}\,4\sin^2(\pi/L)$.

\section{Perturbed configurations}

The same projected Lanczos code reports that for a small random perturbation around vacuum (e.g.\ amplitude $10^{-2}$), the lowest projected eigenvalue can become slightly negative, on the scale of $10^{-7}$--$10^{-6}$ in the reported runs.

\begin{remark}
This numerical negativity may be:
(i) a real curvature/nonlinearity effect away from the exact vacuum linearization,
(ii) leakage due to using a vacuum-defined transverse projector at non-vacuum field,
or (iii) numerical error / Lanczos truncation error.
No inference is justified without a controlled analytic comparison.
\end{remark}

\section{Hessian scans ``on GPU''}

A separate GPU/Lanczos script scans the smallest eigenvalue of the \emph{full} Hessian (not gauge-projected) along random directions at amplitudes in $[10^{-2}, 2.5\times 10^{-2}]$ and reports negative values of size $\sim 10^{-3}$.
Because gauge directions are not removed, these numbers cannot be interpreted as physical convexity constants.

\section{HOTRG Jacobian sketches}

The directory contains a sketch implementation of a HOTRG-style contraction + SVD truncation map $T\mapsto T'$ and a plan to compute its Jacobian by automatic differentiation (JAX).

\begin{remark}
Differentiating through SVD/truncation is not a mathematically clean operation at singular-value degeneracies; the Jacobian may be ill-defined or discontinuous there.
Any theorem using a Jacobian norm bound for HOTRG must explicitly address this nondifferentiability.
\end{remark}

\end{document}
